\documentclass[12pt]{article}

\usepackage[margin=1in]{geometry} % 1-inch margins
\usepackage{times}                % Times (Times New Roman–like) font
\usepackage{fancyhdr}             % For custom headers/footers
\usepackage{graphicx}             % For including images
\usepackage{titlesec}             % For section formatting
\usepackage{setspace}             % For 1.5 line spacing

\pagestyle{fancy}
\fancyhf{} % Clear default header and footer

% Header on the right side
\fancyhead[R]{Saksham AI: Multilingual Virtual Interviewer}

% Footer on the left side
\fancyfoot[L]{Pillai HOC College of Engineering \& Technology (Autonomous), Rasayani}

% Footer on the right side (Page Number)
\fancyfoot[R]{\thepage}

% Horizontal lines in header & footer
\renewcommand{\headrulewidth}{0.4pt} 
\renewcommand{\footrulewidth}{0.4pt} 

% Paragraph Formatting
\setlength{\parindent}{1.5em} % Indent first line
\setlength{\parskip}{0em}     % No extra space between paragraphs
\onehalfspacing  

\begin{document}

\setcounter{page}{46}

\vspace*{\fill}

\begin{center}
    {\Large \bf Chapter 9} \\[1em]
    {\Large \bf Conclusion and Future Scope}
\end{center}

\vspace*{\fill}

\newpage

% Set section counter to 8 so that it starts as Chapter 9
\setcounter{section}{9}

% SECTION 9: CONCLUSION AND FUTURE SCOPE

\subsection{Conclusion}

In conclusion, Saksham AI: Multilingual Virtual Interviewer has been successfully developed with key features implemented to enhance the candidate experience and streamline the interview preparation and evaluation process. This includes the capability for users to participate in dynamic, AI-driven video interviews with real-time speech-to-text processing, ensuring an interactive and highly responsive assessment environment. Additionally, candidates can communicate seamlessly in their preferred language through the system's robust multilingual translation capabilities, promoting inclusivity and breaking down language barriers during technical assessments. Furthermore, the system allows users to upload their resumes for automated parsing and skill analysis, which intelligently tailors the generated interview questions to their specific professional background and technical expertise. These functionalities have been integrated into the system to cater to both comprehensive candidate practice and objective administrative evaluation, providing users with a complete solution for managing modern, automated interviews effectively. Throughout the development process, extensive requirement gathering and analysis were conducted to understand user needs and the technical complexities of AI and language model integration. This was followed by the design and implementation of the features using appropriate software architecture and methodologies. Rigorous testing and validation were performed across the speech recognition, translation accuracy, and AI scoring algorithms to ensure the reliability, efficiency, and usability of the developed system. Finally, Saksham AI: Multilingual Virtual Interviewer was successfully finalized and rolled out for operational use.
\vspace{10pt}

\subsection{Limitations}

Speech recognition performance in the system may be influenced by external conditions such as background noise, variations in regional accents, differences in speech clarity, and the quality of the user’s microphone, which can occasionally result in inaccurate transcription of spoken responses. Since the platform relies on real-time voice interaction, any distortion or interference in audio input can affect the precision of language processing and subsequent evaluation. Additionally, the system executes several computationally intensive tasks concurrently, including audio acquisition, speech-to-text conversion, multilingual translation, and AI-based response analysis. The integration of these processes may introduce slight processing delays during transitions between questions, thereby affecting the immediacy of system feedback and the overall responsiveness of the live interview experience.
\vspace{5pt}

\subsection{Future Scope}

Speech recognition performance in the system may be influenced by external conditions such as background noise, variations in regional accents, differences in speech clarity, and the quality of the user’s microphone, which can occasionally result in inaccurate transcription of spoken responses. Since the platform relies on real-time voice interaction, any distortion or interference in audio input can affect the precision of language processing and subsequent evaluation. Additionally, the system executes several computationally intensive tasks concurrently, including audio acquisition, speech-to-text conversion, multilingual translation, and AI-based response analysis. The integration of these processes may introduce slight processing delays during transitions between questions, thereby affecting the immediacy of system feedback and the overall responsiveness of the live interview experience.


\vspace{10pt}

\end{document}